

Este trabalho investigou estratégias de adaptação eficiente de LLMs de pequeno porte para a tarefa de geração automática de títulos científicos a partir de resumos, utilizando o modelo Qwen2,5-0,5B-Instruct e um subconjunto de 10.000 amostras do repositório arXiv. Quatro abordagens foram comparadas: \textit{Baseline} zero-shot (M1), ajuste da última camada --- \textit{Linear Probing} (M2), Q-LoRA (M3) e \textit{Full Fine-Tuning} (M4).

Os resultados obtidos por meio da métrica S-BERT e corroborados pela avaliação qualitativa humana permitem responder à questão central do estudo com clareza. O \textit{Full Fine-Tuning} (M4, $\mu = 0,640$) representa o teto de desempenho absoluto, superando estatisticamente o \textit{baseline} e o ajuste da última camada ($p < 0,001$). O método Q-LoRA (M3, $\mu = 0,627$), contudo, alcançou desempenho estatisticamente equivalente ao M4 ($p = 0,109$), com qualidade textual praticamente indistinguível na avaliação humana, o que valida sua eficiência como alternativa PEFT em cenários com restrições computacionais. Em contrapartida, o ajuste isolado da última camada (M2, $\mu = 0,596$) produziu desempenho inferior ao do modelo original sem treinamento, evidenciando que a atualização de um único componente da rede, sem ajuste das representações intermediárias, é insuficiente para esta tarefa.

Outro achado relevante é a maior variância observada nos resultados do Q-LoRA em relação ao \textit{Full Fine-Tuning}, manifestada qualitativamente como inserção esporádica de contexto externo nos títulos gerados. Esse comportamento sugere uma sensibilidade de modelos compactos à quantização de 4 bits, que, ao reduzir a redundância paramétrica disponível, pode comprometer a estabilidade gerativa mesmo quando a média de desempenho permanece competitiva.

Como limitação, ressalta-se que os experimentos foram conduzidos com um único modelo de pequeno porte (0,5B parâmetros) e com uma amostra restrita do conjunto de dados, em razão de restrições computacionais. Trabalhos futuros poderão expandir esta investigação avaliando modelos de maior escala (e.g., 1,5B e 7B parâmetros) e diferentes níveis de quantização (4-bit vs.\ 8-bit), a fim de determinar o limiar de compressão aceitável e verificar se a instabilidade observada no M3 é um fenômeno específico de arquiteturas compactas. A literatura recente aponta esta como uma direção promissora: Egashira et al.~\cite{egashira_exploiting_nodate} e Shi e Ding~\cite{shi_systematic_2025} caracterizam sistematicamente os impactos da quantização sobre desempenho, consumo energético e consistência de geração em LLMs, reforçando a necessidade de estudos controlados sobre o tema.